\documentclass[10pt,a4paper]{amsart}
\usepackage{amsmath,amssymb,amsthm,graphicx}
% Retained preamble for the Stacks Project chapter snapshot in this directory.
% The chapter's own preamble.tex is a copy of the shared Stacks Project preamble; the block
% below is sliced MECHANICALLY, line for line, from it. Nothing was added, removed or
% redefined; the wrapper class, the xr-hyper/hyperref wiring and the externaldocument list
% are omitted because the assembler supplies the document class.
%
% lines taken from preamble.tex: 3, 5, 9, 11, 12, 156, 157, 158, ... (51 lines total)
%
% package coverage restored from the chapter's own preamble (2026-10-13): xy, xy-2cell

\usepackage[all]{xy}
\xyoption{2cell}
\UseAllTwocells

\newenvironment{reference}{\comment}{\endcomment}
\newenvironment{slogan}{\comment}{\endcomment}
\newenvironment{history}{\comment}{\endcomment}
\theoremstyle{plain}
\newtheorem{theorem}[subsection]{Theorem}
\newtheorem{proposition}[subsection]{Proposition}
\newtheorem{lemma}[subsection]{Lemma}
\theoremstyle{definition}
\newtheorem{definition}[subsection]{Definition}
\newtheorem{example}[subsection]{Example}
\newtheorem{exercise}[subsection]{Exercise}
\newtheorem{situation}[subsection]{Situation}
\theoremstyle{remark}
\newtheorem{remark}[subsection]{Remark}
\newtheorem{remarks}[subsection]{Remarks}
\numberwithin{equation}{subsection}
\def\lim{\mathop{\mathrm{lim}}\nolimits}
\def\colim{\mathop{\mathrm{colim}}\nolimits}
\def\Spec{\mathop{\mathrm{Spec}}}
\def\Hom{\mathop{\mathrm{Hom}}\nolimits}
\def\Ext{\mathop{\mathrm{Ext}}\nolimits}
\def\SheafHom{\mathop{\mathcal{H}\!\mathit{om}}\nolimits}
\def\SheafExt{\mathop{\mathcal{E}\!\mathit{xt}}\nolimits}
\def\Sch{\mathit{Sch}}
\def\Mor{\mathop{\mathrm{Mor}}\nolimits}
\def\Ob{\mathop{\mathrm{Ob}}\nolimits}
\def\Sh{\mathop{\mathit{Sh}}\nolimits}
\def\NL{\mathop{N\!L}\nolimits}
\def\CH{\mathop{\mathrm{CH}}\nolimits}
\def\proetale{{pro\text{-}\acute{e}tale}}
\def\etale{{\acute{e}tale}}
\def\QCoh{\mathit{QCoh}}
\def\Ker{\mathop{\mathrm{Ker}}}
\def\Im{\mathop{\mathrm{Im}}}
\def\Coker{\mathop{\mathrm{Coker}}}
\def\Coim{\mathop{\mathrm{Coim}}}
\DeclareMathSymbol{\boxtimes}{\mathbin}{AMSa}{"02}
\def\QCohstack{\mathcal{QC}\!\mathit{oh}}
\def\Cohstack{\mathcal{C}\!\mathit{oh}}
\def\Spacesstack{\mathcal{S}\!\mathit{paces}}
\def\Quotfunctor{\mathrm{Quot}}
\def\Hilbfunctor{\mathrm{Hilb}}
\def\Curvesstack{\mathcal{C}\!\mathit{urves}}
\def\Polarizedstack{\mathcal{P}\!\mathit{olarized}}
\def\Complexesstack{\mathcal{C}\!\mathit{omplexes}}
\def\Pic{\mathop{\mathrm{Pic}}\nolimits}
\def\Picardstack{\mathcal{P}\!\mathit{ic}}
\def\Picardfunctor{\mathrm{Pic}}
\def\Deformationcategory{\mathcal{D}\!\mathit{ef}}
\usepackage{cleveref}
\crefname{originallemma}{original lemma}{original lemmas}
\Crefname{originallemma}{Original lemma}{Original lemmas}
\makeatletter
\newlabel{algebra-lemma-localization-colimit}{{00CR}{source}{}{sourcepointer-1}{}}
\newlabel{algebra-lemma-localization-colimit@cref}{{[originallemma][0][]00CR}{[1][0][]source}{}{}{}}
\makeatother
\crefname{originallemma}{original lemma}{original lemmas}
\Crefname{originallemma}{Original lemma}{Original lemmas}
\makeatletter
\newlabel{algebra-lemma-ring-colimit-fp}{{00QN}{source}{}{sourcepointer-2}{}}
\newlabel{algebra-lemma-ring-colimit-fp@cref}{{[originallemma][0][]00QN}{[1][0][]source}{}{}{}}
\makeatother
\crefname{originallemma}{original lemma}{original lemmas}
\Crefname{originallemma}{Original lemma}{Original lemmas}
\makeatletter
\newlabel{formal-defos-lemma-surjective-cotangent-space}{{06GZ}{source}{}{sourcepointer-3}{}}
\newlabel{formal-defos-lemma-surjective-cotangent-space@cref}{{[originallemma][0][]06GZ}{[1][0][]source}{}{}{}}
\makeatother
\crefname{originallemma}{original lemma}{original lemmas}
\Crefname{originallemma}{Original lemma}{Original lemmas}
\makeatletter
\newlabel{more-algebra-lemma-approximate-complex-graded}{{07VF}{source}{}{sourcepointer-4}{}}
\newlabel{more-algebra-lemma-approximate-complex-graded@cref}{{[originallemma][0][]07VF}{[1][0][]source}{}{}{}}
\makeatother
\crefname{originallemma}{original lemma}{original lemmas}
\Crefname{originallemma}{Original lemma}{Original lemmas}
\makeatletter
\newlabel{more-algebra-proposition-finite-type-over-G-ring}{{07PV}{source}{}{sourcepointer-5}{}}
\newlabel{more-algebra-proposition-finite-type-over-G-ring@cref}{{[originallemma][0][]07PV}{[1][0][]source}{}{}{}}
\makeatother
\crefname{originallemma}{original lemma}{original lemmas}
\Crefname{originallemma}{Original lemma}{Original lemmas}
\makeatletter
\newlabel{more-algebra-section-artin-rees}{{07VD}{source}{}{sourcepointer-6}{}}
\newlabel{more-algebra-section-artin-rees@cref}{{[originallemma][0][]07VD}{[1][0][]source}{}{}{}}
\makeatother
\crefname{originallemma}{original lemma}{original lemmas}
\Crefname{originallemma}{Original lemma}{Original lemmas}
\makeatletter
\newlabel{morphisms-lemma-point-finite-type}{{01TA}{source}{}{sourcepointer-7}{}}
\newlabel{morphisms-lemma-point-finite-type@cref}{{[originallemma][0][]01TA}{[1][0][]source}{}{}{}}
\makeatother
\crefname{originallemma}{original lemma}{original lemmas}
\Crefname{originallemma}{Original lemma}{Original lemmas}
\makeatletter
\newlabel{smoothing-theorem-approximation-property-variant}{{07QZ}{source}{}{sourcepointer-8}{}}
\newlabel{smoothing-theorem-approximation-property-variant@cref}{{[originallemma][0][]07QZ}{[1][0][]source}{}{}{}}
\makeatother
\title{Artin approximation of an object over a complete local ring}
\author{The Stacks Project Authors}
\date{stacks-project:artin.tex:local-snapshot-20261012; GNU FDL 1.2 or later}
\begin{document}
\maketitle
\noindent\textit{Editorial note.} Statement span is lines 1487-1509 of sources/stacks-artin/pinned-original-source.tex; proof span is lines 1511-1649, the complete author proof (137 body lines).\par\medskip
\noindent\textit{Editorial note.} Navigation stubs: 8 cited-but-not-included labels are supplied as non-load-bearing pointers with their pinned tags where the snapshot tag index has them.\par\medskip
\noindent\textit{Editorial note.} extra\_packages is empty on purpose: the retained preamble already loads xy with options and a second declaration clashes in the assembler wrapper.\par\medskip
\noindent\textit{Editorial note.} Text is retained verbatim, including the author's own typos and placeholders.\par\medskip
\noindent\textit{Editorial note.} Queue index 1 of queue\_artin.json (pb=5989, stmt 1487-1509, proof 1511-1649).\par\medskip
\noindent\textit{Editorial note.} stacks-artin wave ar01 (2026-10-14): fresh adjudication accept per lane-s3/artin\_judgments.md; pool per Chief 10:5x direct dispatch (idle-first over S12's 01:5x note, SMC precedent); artin.tex 99 rows/67 with proof judged 46 accept / 21 reject / 5 AGP-overlap avoidance rows; patch-sd-envs contract spans scanned.\par\medskip
\noindent\textit{Editorial note.} Cross-chapter same-label warning (NOT an occupancy conflict, Chief ruling 2026-10-13 evening): label 'lemma-approximate' also exists in stacks-restricted, stacks-snapshot-20261011 — a different proposition in a different chapter file.\par\medskip
\noindent\textit{Allowed standard prerequisites (editorial).}\begin{itemize}
\item Explicit prerequisite boundary. The cited inputs are prerequisites; the argument in the retained span is not replaced by any cited result.
\end{itemize}\medskip
\ifdefined\hypertarget\hypertarget{sourcepointer-1}{}\fi
\noindent\textit{Source reference.} Original lemma 00CR: sources/stacks-artin/pinned-original-source.tex, source lines 1649--1649. Cited by the retained statement or proof but not part of the collected spans; supplied only as a navigation pointer. The argument does not depend on its content.\par\medskip
\ifdefined\hypertarget\hypertarget{sourcepointer-2}{}\fi
\noindent\textit{Source reference.} Original lemma 00QN: sources/stacks-artin/pinned-original-source.tex, source lines 1649--1649. Cited by the retained statement or proof but not part of the collected spans; supplied only as a navigation pointer. The argument does not depend on its content.\par\medskip
\ifdefined\hypertarget\hypertarget{sourcepointer-3}{}\fi
\noindent\textit{Source reference.} Original lemma 06GZ: sources/stacks-artin/pinned-original-source.tex, source lines 1649--1649. Cited by the retained statement or proof but not part of the collected spans; supplied only as a navigation pointer. The argument does not depend on its content.\par\medskip
\ifdefined\hypertarget\hypertarget{sourcepointer-4}{}\fi
\noindent\textit{Source reference.} Original lemma 07VF: sources/stacks-artin/pinned-original-source.tex, source lines 1649--1649. Cited by the retained statement or proof but not part of the collected spans; supplied only as a navigation pointer. The argument does not depend on its content.\par\medskip
\ifdefined\hypertarget\hypertarget{sourcepointer-5}{}\fi
\noindent\textit{Source reference.} Original lemma 07PV: sources/stacks-artin/pinned-original-source.tex, source lines 1649--1649. Cited by the retained statement or proof but not part of the collected spans; supplied only as a navigation pointer. The argument does not depend on its content.\par\medskip
\ifdefined\hypertarget\hypertarget{sourcepointer-6}{}\fi
\noindent\textit{Source reference.} Original lemma 07VD: sources/stacks-artin/pinned-original-source.tex, source lines 1649--1649. Cited by the retained statement or proof but not part of the collected spans; supplied only as a navigation pointer. The argument does not depend on its content.\par\medskip
\ifdefined\hypertarget\hypertarget{sourcepointer-7}{}\fi
\noindent\textit{Source reference.} Original lemma 01TA: sources/stacks-artin/pinned-original-source.tex, source lines 1649--1649. Cited by the retained statement or proof but not part of the collected spans; supplied only as a navigation pointer. The argument does not depend on its content.\par\medskip
\ifdefined\hypertarget\hypertarget{sourcepointer-8}{}\fi
\noindent\textit{Source reference.} Original lemma 07QZ: sources/stacks-artin/pinned-original-source.tex, source lines 1649--1649. Cited by the retained statement or proof but not part of the collected spans; supplied only as a navigation pointer. The argument does not depend on its content.\par\medskip
\begin{lemma}
\label{lemma-approximate}
Let $S$ be a locally Noetherian scheme. Let
$p : \mathcal{X} \to (\Sch/S)_{fppf}$ be a category
fibred in groupoids. Let $x$ be an object of
$\mathcal{X}$ lying over $\Spec(R)$ where $R$ is a Noetherian complete
local ring with residue field $k$ of finite type over $S$. Let $s \in S$
be the image of $\Spec(k) \to S$. Assume that (a) $\mathcal{O}_{S, s}$ is
a G-ring and (b) $p$ is limit preserving on objects. Then for every
integer $N \geq 1$ there exist
\begin{enumerate}
\item a finite type $S$-algebra $A$,
\item a maximal ideal $\mathfrak m_A \subset A$,
\item an object $x_A$ of $\mathcal{X}$ over $\Spec(A)$,
\item an $S$-isomorphism $R/\mathfrak m_R^N \cong A/\mathfrak m_A^N$,
\item an isomorphism
$x|_{\Spec(R/\mathfrak m_R^N)} \cong x_A|_{\Spec(A/\mathfrak m_A^N)}$
compatible with (4), and
\item an isomorphism
$\text{Gr}_{\mathfrak m_R}(R) \cong \text{Gr}_{\mathfrak m_A}(A)$
of graded $k$-algebras.
\end{enumerate}
\end{lemma}

\begin{proof}
Choose an affine open $\Spec(\Lambda) \subset S$ such that $k$ is a finite
$\Lambda$-algebra, see
Morphisms, Lemma \ref{morphisms-lemma-point-finite-type}.
We may and do replace $S$ by $\Spec(\Lambda)$.

\medskip\noindent
We may write $R$ as a directed colimit $R = \colim C_j$ where each
$C_j$ is a finite type $\Lambda$-algebra (see
Algebra, Lemma \ref{algebra-lemma-ring-colimit-fp}).
By assumption (b) the object $x$ is isomorphic to the restriction of
an object over one of the $C_j$. Hence we may choose a finite type
$\Lambda$-algebra $C$, a $\Lambda$-algebra map $C \to R$, and an object
$x_C$ of $\mathcal{X}$ over $\Spec(C)$ such that $x = x_C|_{\Spec(R)}$.
The choice of $C$ is a bookkeeping device and could be avoided.
For later use, let us write $C = \Lambda[y_1, \ldots, y_u]/(f_1, \ldots, f_v)$
and we denote $\overline{a}_i \in R$ the image of $y_i$ under the
map $C \to R$. Set $\mathfrak m_C = C \cap \mathfrak m_R$.

\medskip\noindent
Choose a $\Lambda$-algebra surjection $\Lambda[x_1, \ldots, x_s] \to k$
and denote $\mathfrak m'$ the kernel.
By the universal property of polynomial rings we may lift this
to a $\Lambda$-algebra map $\Lambda[x_1, \ldots, x_s] \to R$.
We add some variables (i.e., we increase $s$ a bit) mapping to generators
of $\mathfrak m_R$. Having done this we see that
$\Lambda[x_1, \ldots, x_s] \to R/\mathfrak m_R^2$ is surjective.
Then we see that
\begin{equation}
\label{equation-surjection}
P = \Lambda[x_1, \ldots, x_s]_{\mathfrak m'}^\wedge \longrightarrow R
\end{equation}
is a surjective map of Noetherian complete local rings, see for example
Formal Deformation Theory, Lemma
\ref{formal-defos-lemma-surjective-cotangent-space}.

\medskip\noindent
Choose lifts $a_i \in P$ of $\overline{a}_i$ we found above.
Choose generators $b_1, \ldots, b_r \in P$ for the kernel of
(\ref{equation-surjection}).
Choose $c_{ji} \in P$ such that
$$
f_j(a_1, \ldots, a_u) = \sum c_{ji} b_i
$$
in $P$ which is possible by the choices made so far. Choose generators
$$
k_1, \ldots, k_t \in
\Ker(P^{\oplus r} \xrightarrow{(b_1, \ldots, b_r)} P)
$$
and write $k_i = (k_{i1}, \ldots, k_{ir})$ and $K = (k_{ij})$
so that
$$
P^{\oplus t} \xrightarrow{K}
P^{\oplus r} \xrightarrow{(b_1, \ldots, b_r)}
P \to R \to 0
$$
is an exact sequence of $P$-modules. In particular we have
$\sum k_{ij} b_j = 0$. After possibly increasing $N$ we may
assume $N - 1$ works in the Artin-Rees lemma for the first two maps of this
exact sequence (see More on Algebra, Section
\ref{more-algebra-section-artin-rees} for terminology).

\medskip\noindent
By assumption $\mathcal{O}_{S, s} = \Lambda_{\Lambda \cap \mathfrak m'}$ is
a G-ring. Hence by More on Algebra, Proposition
\ref{more-algebra-proposition-finite-type-over-G-ring}
the ring $\Lambda[x_1, \ldots, x_s]_{\mathfrak m'}$ is a $G$-ring.
Hence by Smoothing Ring Maps, Theorem
\ref{smoothing-theorem-approximation-property-variant}
there exist an \'etale ring map
$$
\Lambda[x_1, \ldots, x_s]_{\mathfrak m'} \to B,
$$
a maximal ideal $\mathfrak m_B$ of $B$ lying over $\mathfrak m'$, and
elements $a'_i, b'_i, c'_{ij}, k'_{ij} \in B'$ such that
\begin{enumerate}
\item $\kappa(\mathfrak m') = \kappa(\mathfrak m_B)$ which implies
that $\Lambda[x_1, \ldots, x_s]_{\mathfrak m'} \subset B_{\mathfrak m_B}
\subset P$ and $P$ is identified with the completion of $B$ at
$\mathfrak m_B$, see remark preceding Smoothing Ring Maps, Theorem
\ref{smoothing-theorem-approximation-property-variant},
\item $a_i - a'_i, b_i - b'_i, c_{ij} - c'_{ij}, k_{ij} - k'_{ij} \in
(\mathfrak m')^N P$, and
\item $f_j(a'_1, \ldots, a'_u) = \sum c'_{ji} b'_i$ and $\sum k'_{ij}b'_j = 0$.
\end{enumerate}
Set $A = B/(b'_1, \ldots, b'_r)$ and denote $\mathfrak m_A$ the
image of $\mathfrak m_B$ in $A$. (Note that $A$ is essentially of finite
type over $\Lambda$; at the end of the proof we will show how to obtain
an $A$ which is of finite type over $\Lambda$.) There is a ring map
$C \to A$ sending $y_i \mapsto a'_i$ because the $a'_i$ satisfy
the desired equations modulo $(b'_1, \ldots, b'_r)$.
Note that $A/\mathfrak m_A^N = R/\mathfrak m_R^N$ as quotients of
$P = B^\wedge$ by property (2) above. Set $x_A = x_C|_{\Spec(A)}$.
Since the maps
$$
C \to A \to A/\mathfrak m_A^N \cong R/\mathfrak m_R^N
\quad\text{and}\quad
C \to R \to R/\mathfrak m_R^N
$$
are equal we see that $x_A$ and $x$ agree modulo $\mathfrak m_R^N$
via the isomorphism $A/\mathfrak m_A^N = R/\mathfrak m_R^N$. At this
point we have shown properties (1) -- (5) of the statement of the lemma.
To see (6) note that
$$
P^{\oplus t} \xrightarrow{K}
P^{\oplus r} \xrightarrow{(b_1, \ldots, b_r)}
P
\quad\text{and}\quad
P^{\oplus t} \xrightarrow{K'}
P^{\oplus r} \xrightarrow{(b'_1, \ldots, b'_r)}
P
$$
are two complexes of $P$-modules which are congruent modulo
$(\mathfrak m')^N$ with the first one being exact. By our choice of $N$
above we see from
More on Algebra, Lemma \ref{more-algebra-lemma-approximate-complex-graded}
that $R = P/(b_1, \ldots, b_r)$ and
$P/(b'_1, \ldots, b'_r) = B^\wedge/(b'_1, \ldots, b'_r) = A^\wedge$
have isomorphic associated graded algebras, which is what we wanted to show.

\medskip\noindent
This last paragraph of the proof serves to clean up the issue that $A$ is
essentially of finite type over $S$ and not yet of finite type.
The construction above gives $A = B/(b'_1, \ldots, b'_r)$ and
$\mathfrak m_A \subset A$ with $B$ \'etale over
$\Lambda[x_1, \ldots, x_s]_{\mathfrak m'}$. Hence $A$ is of finite
type over the Noetherian ring $\Lambda[x_1, \ldots, x_s]_{\mathfrak m'}$.
Thus we can write $A = (A_0)_{\mathfrak m'}$ for some finite type
$\Lambda[x_1, \ldots, x_n]$ algebra $A_0$. Then
$A = \colim (A_0)_f$ where
$f \in \Lambda[x_1, \ldots, x_n] \setminus \mathfrak m'$, see
Algebra, Lemma \ref{algebra-lemma-localization-colimit}.
Because $p : \mathcal{X} \to (\Sch/S)_{fppf}$ is limit preserving on
objects, we see that
$x_A$ comes from some object $x_{(A_0)_f}$ over $\Spec((A_0)_f)$ for
an $f$ as above. After replacing $A$ by $(A_0)_f$ and $x_A$ by
$x_{(A_0)_f}$ and $\mathfrak m_A$ by $(A_0)_f \cap \mathfrak m_A$
the proof is finished.
\end{proof}

\end{document}
